\documentclass[10pt]{article}

\usepackage[a4paper,margin=1.8cm]{geometry}
\usepackage{amsmath,amssymb,mathtools}
\usepackage{booktabs,tabularx,array}
\usepackage{microtype}
\usepackage[hidelinks]{hyperref}
\usepackage{enumitem}

\setlength{\parindent}{0pt}
\setlength{\parskip}{0.45em}
\setlist{nosep,leftmargin=*}
\newcommand{\ind}{\mathbf{1}}
\newcommand{\dd}{\mathrm{d}}
\newcommand{\e}{\mathrm{e}}
\newcommand{\R}{\mathbb{R}}
\newcommand{\T}{\mathcal{T}}
\newcommand{\D}{\mathcal{D}}
\newcommand{\eps}{\varepsilon}
\newcommand{\parone}{\eta_1}
\newcommand{\partwo}{\eta_2}
\newcommand{\parthree}{\eta_3}

\title{Analytic derivatives for the BB and Tawn bivariate copulas in \texttt{vinecopulib}}
\author{}
\date{11 July 2026}

\begin{document}
\maketitle
\vspace{-2.2em}

\paragraph{Scope.}
This note gives exact, implementation-ready formulas for the unrotated BB1, BB6,
BB7, BB8, and Tawn copulas in the parameterizations used by
\texttt{vinecopulib}.  It covers the copula $C$, the two library h-functions
$h_1=C_u$ and $h_2=C_v$, the density $c=C_{uv}$, and every first- and
second-order derivative of $c,h_1,h_2$ with respect to their arguments and
parameters.  Thus it covers the selectors accepted by
\texttt{pdf\_deriv}, \texttt{pdf\_deriv2},
\texttt{hfunc1\_deriv}, \texttt{hfunc1\_deriv2},
\texttt{hfunc2\_deriv}, and \texttt{hfunc2\_deriv2}; derivatives of
$\log c$ follow at no additional cost.  Rotations are not repeated here,
because the public \texttt{Bicop} facade already applies their chain rules.

\section{Coordinates and derivative selectors}

Let
$ x=(u,v,\eta_1,\ldots,\eta_p) $ and write $D_i=\partial/\partial x_i$.
The correspondence with the library is
\texttt{u1}$\leftrightarrow u$, \texttt{u2}$\leftrightarrow v$, and
\texttt{par}$r\leftrightarrow\eta_r$.  For any positive scalar function $f$,
write
\begin{align*}
D_i\log f &= \frac{D_i f}{f},
&
D_{ij}\log f &= \frac{D_{ij}f}{f}-\frac{D_i f\,D_j f}{f^2},\\
D_i f &= f\,D_i\log f,
&
D_{ij}f &= f\{D_{ij}\log f+(D_i\log f)(D_j\log f)\}.
\end{align*}
These identities are used repeatedly below and are numerically preferable to
fully expanded expressions.

\section{BB1, BB6, BB7, and BB8}

\subsection{Generators, copulas, h-functions, and densities}

All four families are strict Archimedean copulas on the interior of their
parameter spaces.  With decreasing generator $\phi$ and inverse $q=\phi^{-1}$,
set
\begin{align*}
s &= \phi(u;\eta)+\phi(v;\eta),
&
w &= q(s;\eta)=C(u,v;\eta).
\end{align*}
The parameterizations in \texttt{vinecopulib} are as follows.  The ``code
range'' records the finite bounds currently used by the implementation; the
formulas below concern the mathematical interior, with boundary cases obtained
by continuity.

\renewcommand{\arraystretch}{1.25}
\begin{center}
\small
\begin{tabularx}{\textwidth}{@{}l>{\raggedright\arraybackslash}p{2.3cm}X X@{}}
\toprule
Family & Parameters / code range & $\phi(z)$ & $q(s)$ \\
\midrule
BB1
& $\theta>0,\ \delta\ge1$; $[0,7]\times[1,7]$
& $(z^{-\theta}-1)^\delta$
& $(1+s^{1/\delta})^{-1/\theta}$ \\
BB6
& $\theta\ge1,\ \delta\ge1$; $[1,6]\times[1,8]$
& $\{-\log[1-(1-z)^\theta]\}^\delta$
& $1-\{1-\exp(-s^{1/\delta})\}^{1/\theta}$ \\
BB7
& $\theta\ge1,\ \delta>0$; $[1,6]\times[0.01,25]$
& $[1-(1-z)^\theta]^{-\delta}-1$
& $1-\{1-(1+s)^{-1/\delta}\}^{1/\theta}$ \\
BB8
& $\theta\ge1,\ 0<\delta\le1$; $[1,8]\times[10^{-4},1]$
& $-\log\!\left\{\dfrac{1-(1-\delta z)^\theta}{1-(1-\delta)^\theta}\right\}$
& $\dfrac{1-\{1-\kappa\exp(-s)\}^{1/\theta}}{\delta}$,
$\ \kappa=1-(1-\delta)^\theta$ \\
\bottomrule
\end{tabularx}
\end{center}

Differentiating $\phi(w)=\phi(u)+\phi(v)$ gives the exact formulas used by the
Archimedean implementation:
\begin{align*}
h_1(u,v;\eta) &= C_u=\frac{\phi_z(u)}{\phi_z(w)},\\
h_2(u,v;\eta) &= C_v=\frac{\phi_z(v)}{\phi_z(w)},\\
c(u,v;\eta) &= C_{uv}
=-\frac{\phi_z(u)\phi_z(v)\phi_{zz}(w)}{\phi_z(w)^3}.
\end{align*}

\subsection{A universal exact derivative calculus}

For a scalar $f(z,\eta)$ evaluated at a scalar map $z=z(x)$, define
$e_i^a=\ind\{x_i=\eta_a\}$ and
\begin{align*}
\T_i[f;z]
&=f_z z_i+\sum_a f_a e_i^a,\\
\T_{ij}[f;z]
&=f_{zz}z_i z_j+f_z z_{ij}
 +\sum_a f_{za}(z_i e_j^a+z_j e_i^a)
 +\sum_{a,b}f_{ab}e_i^a e_j^b.
\end{align*}
For $z=u$ and $z=v$,
\begin{align*}
u_i&=\ind\{x_i=u\},&u_{ij}&=0,
&v_i&=\ind\{x_i=v\},&v_{ij}&=0.
\end{align*}
For $z=w=C(u,v;\eta)$, implicit differentiation gives
\begin{align*}
w_i
&=\frac{\T_i[\phi;u]+\T_i[\phi;v]-\sum_a\phi_a(w)e_i^a}
        {\phi_z(w)},\\
w_{ij}
&=\frac{\T_{ij}[\phi;u]+\T_{ij}[\phi;v]-R_{ij}^{(w)}}
        {\phi_z(w)},\\
R_{ij}^{(w)}
&=\phi_{zz}(w)w_iw_j
 +\sum_a\phi_{za}(w)(w_i e_j^a+w_j e_i^a)
 +\sum_{a,b}\phi_{ab}(w)e_i^a e_j^b.
\end{align*}
Hence $C_i=w_i$ and $C_{ij}=w_{ij}$.

Now define the two scalar log-building blocks
\begin{align*}
G(z,\eta)&=\log[-\phi_z(z;\eta)],
&
R(z,\eta)&=\log \phi_{zz}(z;\eta).
\end{align*}
Then
\begin{align*}
H_1&=\log h_1=G(u)-G(w),\\
H_2&=\log h_2=G(v)-G(w),\\
L&=\log c=G(u)+G(v)+R(w)-3G(w).
\end{align*}
All requested first- and second-order derivatives are therefore
\begin{align*}
(H_1)_i
&=\T_i[G;u]-\T_i[G;w],\\
(H_1)_{ij}
&=\T_{ij}[G;u]-\T_{ij}[G;w],\\
(H_2)_i
&=\T_i[G;v]-\T_i[G;w],\\
(H_2)_{ij}
&=\T_{ij}[G;v]-\T_{ij}[G;w],\\
L_i
&=\T_i[G;u]+\T_i[G;v]+\T_i[R;w]-3\T_i[G;w],\\
L_{ij}
&=\T_{ij}[G;u]+\T_{ij}[G;v]+\T_{ij}[R;w]-3\T_{ij}[G;w],\\
D_i h_k
&=h_k(H_k)_i,\\
D_{ij}h_k
&=h_k\{(H_k)_{ij}+(H_k)_i(H_k)_j\},\qquad k\in\{1,2\},\\
D_i c
&=cL_i,\\
D_{ij}c
&=c\{L_{ij}+L_iL_j\}.
\end{align*}
In particular, the analytic log-density leaves are simply
$D_i\log c=L_i$ and $D_{ij}\log c=L_{ij}$.

The partial derivatives of $G$ and $R$ required by $\T_i$ and $\T_{ij}$ are
obtained from derivatives of $\phi$.  With $g=\phi_z$ and $r=\phi_{zz}$,
\begin{align*}
G_z&=\frac{r}{g},
&G_{zz}&=\frac{\phi_{zzz}}{g}-\left(\frac{r}{g}\right)^2,\\
G_a&=\frac{\phi_{za}}{g},
&G_{za}&=\frac{\phi_{zza}}{g}-\frac{r\phi_{za}}{g^2},\\
G_{ab}&=\frac{\phi_{zab}}{g}-\frac{\phi_{za}\phi_{zb}}{g^2},\\[0.2em]
R_z&=\frac{\phi_{zzz}}{r},
&R_{zz}&=\frac{\phi_{zzzz}}{r}-\left(\frac{\phi_{zzz}}{r}\right)^2,\\
R_a&=\frac{\phi_{zza}}{r},
&R_{za}&=\frac{\phi_{zzza}}{r}-\frac{\phi_{zzz}\phi_{zza}}{r^2},\\
R_{ab}&=\frac{\phi_{zzab}}{r}-\frac{\phi_{zza}\phi_{zzb}}{r^2}.
\end{align*}
Thus only generator derivatives $\partial_z^k\partial_\eta^\alpha\phi$ with
$k\le4$ and $|\alpha|\le2$ are needed.

\subsection{Exact generator jets without expanded expressions}

The required generator derivatives are most compactly evaluated as finite
multivariate jets.  For a multi-index $\alpha$, let $f_\alpha=\partial^\alpha f$,
$\binom{\alpha}{\beta}=\prod_j\binom{\alpha_j}{\beta_j}$, and $e_j$ denote a
unit multi-index.  If $f=\exp(g)$, then
\begin{align*}
f_{\alpha+e_j}
=\sum_{\beta\le\alpha}\binom{\alpha}{\beta}
 f_\beta g_{\alpha-\beta+e_j}.
\end{align*}
If $g=\log f$, the same identity can be solved for the highest derivative:
\begin{align*}
g_{\alpha+e_j}
=\frac{f_{\alpha+e_j}
 -\displaystyle\sum_{0<\beta\le\alpha}\binom{\alpha}{\beta}
 f_\beta g_{\alpha-\beta+e_j}}{f}.
\end{align*}
Products use the multivariate Leibniz rule
\begin{align*}
(fg)_\alpha
=\sum_{\beta\le\alpha}\binom{\alpha}{\beta}f_\beta g_{\alpha-\beta}.
\end{align*}
Applying these identities to the following log-domain representations gives
all generator partials exactly, without finite differences or unmanageably long
CAS expansions:
\begin{align*}
\text{BB1:}\quad
&\phi=\exp(q),
&q&=\delta\log\{\exp(-\theta\log z)-1\},\\
\text{BB6:}\quad
&\phi=\exp(q),
&q&=\delta\log\{-\log[1-\exp\{\theta\log(1-z)\}]\},\\
\text{BB7:}\quad
&\phi=\exp(q)-1,
&q&=-\delta\log[1-\exp\{\theta\log(1-z)\}],\\
\text{BB8:}\quad
&\phi=-\log[1-\exp\{\theta\log(1-\delta z)\}]
 +\log[1-\exp\{\theta\log(1-\delta)\}].
\end{align*}
For the pure $z$ derivatives of an exponential, the first four cases reduce to
\begin{align*}
(\e^q)'&=\e^q q',\\
(\e^q)''&=\e^q(q''+q'^2),\\
(\e^q)'''&=\e^q(q'''+3q'q''+q'^3),\\
(\e^q)''''&=\e^q(q''''+4q'q'''+3q''^2+6q'^2q''+q'^4).
\end{align*}
The multivariate recurrence is preferable in code because it simultaneously
produces the mixed parameter derivatives needed above.

\section{Tawn copula}

\subsection{Pickands function, copula, h-functions, and density}

The implementation uses parameters
$\eta=(\psi_1,\psi_2,\theta)$ with
$0\le\psi_1,\psi_2\le1$ and $\theta\ge1$; the code bounds $\theta$ by $60$.
Let
\begin{align*}
a&=-\log u,
&b&=-\log v,
&r&=a+b,
&t&=\frac{b}{r},\\
S&=(\psi_2t)^\theta+[\psi_1(1-t)]^\theta,\\
A(t)&=(1-\psi_1)(1-t)+(1-\psi_2)t+S^{1/\theta}.
\end{align*}
Write $A_k=\partial_t^k A$.  Define
\begin{align*}
P&=A_0-tA_1,
&Q&=A_0+(1-t)A_1,
&d&=t(1-t),\\
B&=PQ+\frac{dA_2}{r}.
\end{align*}
Then the formulas implemented by the extreme-value base class are
\begin{align*}
C(u,v;\eta)&=\exp[-rA(t)],\\
h_1=C_u&=\frac{CP}{u},\\
h_2=C_v&=\frac{CQ}{v},\\
c=C_{uv}&=\frac{CB}{uv}.
\end{align*}
Equivalently,
\begin{align*}
B
=A^2+(1-2t)AA_1-t(1-t)A_1^2+\frac{t(1-t)}{r}A_2,
\end{align*}
which is the source expression after using $\log(uv)=-r$.

For reference, let
\begin{align*}
x&=\psi_2t,
&y&=\psi_1(1-t),\\
M&=\psi_2x^{\theta-1}-\psi_1y^{\theta-1},\\
N&=\psi_2^2x^{\theta-2}+\psi_1^2y^{\theta-2},\\
O&=\psi_2^3x^{\theta-3}-\psi_1^3y^{\theta-3},\\
J&=\psi_2^4x^{\theta-4}+\psi_1^4y^{\theta-4},
&q&=\frac1\theta.
\end{align*}
The argument derivatives needed below are
\begin{align*}
A_1
&=\psi_1-\psi_2+S^{q-1}M,\\
A_2
&=(\theta-1)\{S^{q-1}N-S^{q-2}M^2\},\\
A_3
&=(\theta-1)\{(2\theta-1)S^{q-3}M^3
 -3(\theta-1)S^{q-2}MN
 +(\theta-2)S^{q-1}O\},\\
A_4
&=(\theta-1)\{-(2\theta-1)(3\theta-1)S^{q-4}M^4
 +6(\theta-1)(2\theta-1)S^{q-3}M^2N\\
&\hspace{4.3em}
 -3(\theta-1)^2S^{q-2}N^2
 -4(\theta-1)(\theta-2)S^{q-2}MO
 +(\theta-2)(\theta-3)S^{q-1}J\}.
\end{align*}

\subsection{Coordinate pullback}

For $x_i\in\{u,v,\psi_1,\psi_2,\theta\}$, define
\begin{align*}
a_i&=-\frac{\ind\{x_i=u\}}{u},
&a_{ij}&=\frac{\ind\{x_i=u\}\ind\{x_j=u\}}{u^2},\\
b_i&=-\frac{\ind\{x_i=v\}}{v},
&b_{ij}&=\frac{\ind\{x_i=v\}\ind\{x_j=v\}}{v^2},\\
r_i&=a_i+b_i,
&r_{ij}&=a_{ij}+b_{ij},\\
t_i&=\frac{b_i-tr_i}{r},\\
t_{ij}&=\frac{b_{ij}-t_jr_i-t_ir_j-tr_{ij}}{r}.
\end{align*}
For any scalar $F(r,t,\eta)$, define
\begin{align*}
\D_iF
&=F_r r_i+F_t t_i+\sum_aF_a e_i^a,\\
\D_{ij}F
&=F_{rr}r_ir_j+F_{rt}(r_it_j+t_ir_j)+F_{tt}t_it_j
 +F_r r_{ij}+F_t t_{ij}\\
&\quad
 +\sum_aF_{ra}(r_ie_j^a+r_je_i^a)
 +\sum_aF_{ta}(t_ie_j^a+t_je_i^a)
 +\sum_{a,b}F_{ab}e_i^ae_j^b.
\end{align*}

Set $E=-rA_0$.  The required base derivatives are
\begin{align*}
E_r&=-A_0,&E_t&=-rA_1,&E_a&=-rA_{0,a},\\
E_{rr}&=0,&E_{rt}&=-A_1,&E_{ra}&=-A_{0,a},\\
E_{tt}&=-rA_2,&E_{ta}&=-rA_{1,a},&E_{ab}&=-rA_{0,ab},
\end{align*}
\begin{align*}
P_r&=0,&P_t&=-tA_2,&P_a&=A_{0,a}-tA_{1,a},\\
P_{rr}&=P_{rt}=P_{ra}=0,
&P_{tt}&=-A_2-tA_3,\\
P_{ta}&=-tA_{2,a},
&P_{ab}&=A_{0,ab}-tA_{1,ab},
\end{align*}
\begin{align*}
Q_r&=0,&Q_t&=(1-t)A_2,&Q_a&=A_{0,a}+(1-t)A_{1,a},\\
Q_{rr}&=Q_{rt}=Q_{ra}=0,
&Q_{tt}&=-A_2+(1-t)A_3,\\
Q_{ta}&=(1-t)A_{2,a},
&Q_{ab}&=A_{0,ab}+(1-t)A_{1,ab}.
\end{align*}
Writing $d'=1-2t$ and $d''=-2$, the derivatives of
$B=PQ+dA_2/r$ are
\begin{align*}
B_r
&=-\frac{dA_2}{r^2},\\
B_t
&=P_tQ+PQ_t+\frac{d'A_2+dA_3}{r},\\
B_a
&=P_aQ+PQ_a+\frac{dA_{2,a}}{r},\\
B_{rr}
&=\frac{2dA_2}{r^3},\\
B_{rt}
&=-\frac{d'A_2+dA_3}{r^2},\\
B_{ra}
&=-\frac{dA_{2,a}}{r^2},\\
B_{tt}
&=P_{tt}Q+2P_tQ_t+PQ_{tt}
 +\frac{d''A_2+2d'A_3+dA_4}{r},\\
B_{ta}
&=P_{ta}Q+P_tQ_a+P_aQ_t+PQ_{ta}
 +\frac{d'A_{2,a}+dA_{3,a}}{r},\\
B_{ab}
&=P_{ab}Q+P_aQ_b+P_bQ_a+PQ_{ab}
 +\frac{dA_{2,ab}}{r}.
\end{align*}

Define the log-values
\begin{align*}
\Lambda_1&=\log h_1=E+\log P+a,\\
\Lambda_2&=\log h_2=E+\log Q+b,\\
L&=\log c=E+\log B+a+b.
\end{align*}
For $F\in\{P,Q,B\}$, write
\begin{align*}
\mathcal{L}_i(F)&=\frac{\D_iF}{F},\\
\mathcal{L}_{ij}(F)&=\frac{\D_{ij}F}{F}
 -\frac{(\D_iF)(\D_jF)}{F^2}.
\end{align*}
Then every requested derivative is given by
\begin{align*}
(\Lambda_1)_i
&=\D_iE+\mathcal{L}_i(P)+a_i,\\
(\Lambda_1)_{ij}
&=\D_{ij}E+\mathcal{L}_{ij}(P)+a_{ij},\\
(\Lambda_2)_i
&=\D_iE+\mathcal{L}_i(Q)+b_i,\\
(\Lambda_2)_{ij}
&=\D_{ij}E+\mathcal{L}_{ij}(Q)+b_{ij},\\
L_i
&=\D_iE+\mathcal{L}_i(B)+a_i+b_i,\\
L_{ij}
&=\D_{ij}E+\mathcal{L}_{ij}(B)+a_{ij}+b_{ij},\\
D_i h_k
&=h_k(\Lambda_k)_i,\\
D_{ij}h_k
&=h_k\{(\Lambda_k)_{ij}+(\Lambda_k)_i(\Lambda_k)_j\},\\
D_i c
&=cL_i,\\
D_{ij}c
&=c\{L_{ij}+L_iL_j\}.
\end{align*}
The copula derivatives themselves are
\begin{align*}
D_iC&=C\D_iE,
&D_{ij}C&=C\{\D_{ij}E+(\D_iE)(\D_jE)\}.
\end{align*}

\subsection{Exact parameter derivatives of the Pickands function}

The preceding formulas require
$A_{k,a}$ and $A_{k,ab}$ for $0\le k\le3$ and $0\le k\le2$,
respectively.  They are obtained exactly from
\begin{align*}
X&=\exp\{\theta(\log\psi_2+\log t)\},\\
Y&=\exp\{\theta(\log\psi_1+\log(1-t))\},\\
S&=X+Y,\\
q&=\theta^{-1}\log S,\\
R&=\exp(q),\\
A&=(1-\psi_1)(1-t)+(1-\psi_2)t+R,
\end{align*}
using the multivariate exponential, logarithm, and product recurrences stated
in Section 2.3.  Only indices satisfying
$\alpha_t\le4$ and
$\alpha_{\psi_1}+\alpha_{\psi_2}+\alpha_\theta\le2$ need be retained.
This is an analytic forward jet, not numerical differentiation.  At
$\psi_i=0$, evaluate the continuous limits of
$\psi_i^\theta\log^m\psi_i$ (zero for $\theta>0$) rather than forming
$\log0$.

\section{Implementation remarks}

\begin{itemize}
\item The formulas are for the unrotated leaves.  Existing rotation handling
      supplies argument swaps and signs, while Tawn flipping exchanges
      $\psi_1$ and $\psi_2$.
\item Evaluate the primitive generators in the same stable forms as the current
      source: \texttt{log1p} and \texttt{expm1} for BB6/BB8 and log-domain
      products for small probabilities.
\item At independence or other parameter boundaries, use analytic limits or the
      existing boundary dispatch; interior formulas contain removable
      singularities such as $1/\theta$ for the BB1 limit $\theta\downarrow0$.
\item The identities for $C,h_1,h_2,c$ and the displayed $A_1,\ldots,A_4$
      were numerically checked at interior parameter values against automatic
      differentiation to double precision.
\end{itemize}

\begin{thebibliography}{9}
\bibitem{SchepsmeierStoeber2014}
U.~Schepsmeier and J.~St\"ober.
\newblock Derivatives and Fisher information of bivariate copulas.
\newblock \emph{Statistical Papers}, 55:525--542, 2014.
\newblock \href{https://doi.org/10.1007/s00362-013-0498-x}{doi:10.1007/s00362-013-0498-x}.

\bibitem{SchepsmeierStoeberSupplement}
U.~Schepsmeier and J.~St\"ober.
\newblock Web supplement: Derivatives and Fisher information of bivariate copulas.

\bibitem{Joe2014}
H.~Joe.
\newblock \emph{Dependence Modeling with Copulas}.
\newblock CRC Press, 2014.

\bibitem{vinecopulib}
T.~Nagler and T.~Vatter.
\newblock \texttt{vinecopulib}: a C++ library for vine copula models.
\newblock Source parameterizations inspected on branch \texttt{bicop-derivs},
11 July 2026.
\end{thebibliography}

\end{document}
